\section{Resultados de la implementación}
\begin{indentedsubsection}
A la fecha de elaboración de este informe, el proyecto se encuentra en la fase de cierre de implementación y ``marcha blanca'' de la plataforma Odoo. Esto significa que los módulos configurados ya están disponibles para su uso operativo diario por parte de la gerencia, aunque los resultados a largo plazo sobre el volumen de negocio se evaluarán en los próximos meses. Por ello, los resultados presentados en esta sección documentan los logros inmediatos de configuración y las metas de desempeño (KPIs) proyectadas.
\end{indentedsubsection}

\subsection{Resultados inmediatos de configuración}
\begin{indentedsubsection}
Entre los principales resultados tangibles alcanzados al término de la parametrización del sistema se encuentran:

\begin{itemize}
    \item \textbf{A. Instancia Odoo operativa y centralizada:} despliegue del entorno Odoo Online (SaaS) con configuración bimonetaria (Soles y Dólares), indispensable para el cobro a turistas extranjeros.
    \item \textbf{B. CRM activo (Fin de los cuadernos físicos):} se centralizó la base de datos de pasajeros (PAX) y proveedores, eliminando el riesgo de pérdida de información inherente al uso de libretas de papel y consolidando un embudo de ventas visual.
    \item \textbf{C. Cotizador automatizado y PDFs limpios:} se configuró el catálogo de tarifas base (hoteles, trenes, guías, transportes). El sistema ahora prorratea automáticamente los costos fijos según la cantidad de pasajeros y permite simular porcentajes de ganancia dinámicos (ej. 15\%, 20\%, 30\%). Además, se logró que el PDF final de la cotización oculte el desglose individual del IGV por servicio, mostrando un precio limpio y amigable para el turista internacional, tal como lo solicitó la gerencia.
    \item \textbf{D. Control de pagos fraccionados:} el flujo de ventas ahora permite registrar pagos ``a cuenta'' (como los depósitos por Western Union) y calcula instantáneamente el ``saldo pendiente'', brindando seguridad financiera antes de la ejecución del tour.
    \item \textbf{E. Encuestas post-servicio:} se configuró el formulario digital de retroalimentación para que los turistas califiquen su experiencia al finalizar el viaje, automatizando la recolección de testimonios para potenciar el ``boca a boca''.
\end{itemize}

En términos de madurez operativa, estos resultados no representan únicamente una puesta en marcha técnica, sino la consolidación de una nueva forma de trabajo. La información crítica del negocio deja de estar dispersa y pasa a gestionarse de forma estructurada, con trazabilidad y criterios unificados para ventas, cobros y seguimiento comercial.
\end{indentedsubsection}

\subsection{Resultados preliminares (Marcha blanca)}
\begin{indentedsubsection}
Durante estas primeras semanas de uso, se ha observado un cambio radical en la forma de trabajo de la gerencia. Se ha dejado de lado la dependencia de buscar plantillas antiguas de Word o Excels creados en pandemia. Ahora, la gerencia y su equipo de apoyo pueden visualizar el estado de cada prospecto en una única pantalla, armar paquetes complejos con pocos clics y generar propuestas comerciales formales sin recurrir a calculadoras externas.

Asimismo, la ``marcha blanca'' ha permitido validar que la configuración responde al flujo real de la agencia: desde el registro inicial del lead hasta la emisión de cotización y el control de pagos parciales. Este periodo también ha servido para afinar criterios de uso interno, asegurando que la operación diaria mantenga consistencia aun cuando participan más de un usuario en el proceso comercial.
\end{indentedsubsection}

\subsection{Indicadores y metas de desempeño (KPIs)}
\begin{indentedsubsection}
En coherencia con el Plan de Implementación aprobado por ProInnóvate, la consolidación de este sistema deberá reflejarse en los siguientes Indicadores Clave de Desempeño (KPIs) durante el primer año de uso:
\end{indentedsubsection}

\renewcommand{\arraystretch}{1.35}
\begin{tabularx}{\textwidth}{|>{\raggedright\arraybackslash}p{3.1cm}|>{\raggedright\arraybackslash}p{3.0cm}|>{\raggedright\arraybackslash}p{2.8cm}|X|}
    \hline
    \textbf{KPI (Indicador)} & \textbf{Situación Previa (Manual)} & \textbf{Meta Esperada (con Odoo)} & \textbf{Impacto Esperado} \\ \hline
    Tiempo Promedio de Cotización & 45 - 60 minutos por cliente & < 15 minutos & Reducción drástica de la carga administrativa (cálculos manuales). \\ \hline
    Tasa de Conversión & Sin medición exacta & +10\% a +15\% vs año anterior & Incremento de ventas por respuesta rápida al cliente extranjero. \\ \hline
    Tasa de Error Operativo & \textasciitilde5\% en fechas/nombres & < 1\% & Eliminación de errores de ``copiar y pegar'' en Word. \\ \hline
    Satisfacción del Cliente & Medición informal & Mejora del 15\% & Medición exacta gracias a las Encuestas automatizadas de Odoo. \\ \hline
\end{tabularx}

\begin{indentedsubsection}
Estos KPIs serán monitoreados de forma periódica durante el primer año de adopción para evaluar la consolidación del cambio operativo. El enfoque de seguimiento prioriza no solo la rapidez de cotización, sino también la calidad de atención percibida por el turista y la reducción de errores que afectan la experiencia del cliente final.
\end{indentedsubsection}

\subsection{Impacto esperado en la eficiencia operativa}
\begin{indentedsubsection}
Con la erradicación de la doble digitación y las sumas repetitivas, se proyecta una reducción del 30\% en los tiempos administrativos. Este ahorro de tiempo permitirá que la Gerencia General redirija sus esfuerzos hacia lo que realmente hace crecer el negocio: la venta consultiva, la atención personalizada del turista y el fortalecimiento de sus alianzas con proveedores en destinos como Cusco, Arequipa y Puno. En términos financieros, la optimización operativa proyecta un retorno de inversión (ROI) positivo en un horizonte de 18 a 24 meses.

De manera complementaria, este incremento de eficiencia fortalece la capacidad de la empresa para escalar su operación sin incrementar proporcionalmente su carga administrativa. En consecuencia, Perú Flores podrá atender más oportunidades comerciales con mayor orden interno, preservando la calidad del servicio y mejorando la sostenibilidad de su crecimiento.
\end{indentedsubsection}
